\documentclass[11pt,a4paper]{article}

\usepackage[utf8]{inputenc}
\usepackage[margin=1in]{geometry}
\usepackage{amsmath,amssymb,amsthm}
\usepackage{booktabs}
\usepackage{microtype}
\usepackage{hyperref}
\usepackage{cleveref}
\usepackage{tikz}
\usetikzlibrary{arrows.meta,positioning,shapes.geometric}

% ---------------------------------------------------------------------------
% Theorem Environments
% ---------------------------------------------------------------------------
\newtheorem{theorem}{Theorem}
\newtheorem{lemma}[theorem]{Lemma}
\newtheorem{proposition}[theorem]{Proposition}
\newtheorem{corollary}[theorem]{Corollary}
\theoremstyle{definition}
\newtheorem{definition}[theorem]{Definition}
\newtheorem{example}[theorem]{Example}
\theoremstyle{remark}
\newtheorem{remark}[theorem]{Remark}

% ---------------------------------------------------------------------------
% Metadata
% ---------------------------------------------------------------------------
\title{\textbf{Formalizing the Local-Global Gap in Symbolic World Models:\\The Causal Decoupling Theorem and Topological Divergence Bounds}}
\author{\textbf{Minh-Phuc Tran}\\
Department of Computer Science \& Artificial Intelligence\\
Ho Chi Minh City, Vietnam}
\date{October 2026}

\begin{document}

\maketitle

\begin{abstract}
Symbolic action models (STRIPS/PDDL) provide the foundation for model-based planning, heuristic search, and counterfactual simulation. A pervasive assumption in Action Model Learning (AML) and program-synthesized world models is that high passive transition prediction accuracy ($A_{\text{pred}}$) over observed traces guarantees reliable downstream plan execution. In this work, we prove that this assumption is fundamentally flawed in combinatorial search graphs. We formalize the \emph{Search-Tree Topology Collapse} phenomenon under sparse precondition omissions. We establish: (i) the \emph{No-Edges-Lost Lemma} (\Cref{lem:no_edges_lost}), proving that precondition omissions strictly introduce spurious phantom edges; (ii) the \emph{Search-Tree Phantom Path Divergence Proposition} (\Cref{prop:phantom_divergence}), proving an exponential divergence $d_\triangle = \Omega(b^{D-d})$ while passive accuracy satisfies $A_{\text{pred}} \ge 1 - b^{-d}$; and (iii) the \emph{Causal Decoupling Proposition} (\Cref{prop:independence}), proving that $A_{\text{pred}}$ and Plan Execution Success Rate ($\text{PESR}$) are statistically and functionally independent. These results demonstrate why passive trajectory evaluation is incapable of certifying world models for goal-directed autonomous agents.
\end{abstract}

\tableofcontents
\newpage

% ===========================================================================
\section{Introduction and Foundational Notations}
% ===========================================================================

Let $M^* = \langle \mathcal{F}, \mathcal{A}, \delta^*, s_0, S_g \rangle$ be a ground-truth deterministic transition system specified in relational STRIPS.
\begin{itemize}
    \item $\mathcal{F} = \{f_1, f_2, \dots, f_n\}$ is a finite set of relational Boolean fluents. The state space is factored: $\mathcal{S} = 2^\mathcal{F}$, with cardinality $|\mathcal{S}| = 2^n$.
    \item $\mathcal{A}$ is a finite set of action schemas. Each action $a \in \mathcal{A}$ is defined by a triple:
    \begin{equation}
        a = \langle \text{Pre}(a), \text{Add}(a), \text{Del}(a) \rangle
    \end{equation}
    where $\text{Pre}(a) \subseteq \mathcal{F}$ represents preconditions, and $\text{Add}(a), \text{Del}(a) \subseteq \mathcal{F}$ denote disjoint add and delete effects ($\text{Add}(a) \cap \text{Del}(a) = \emptyset$).
    \item An action $a$ is \emph{applicable} in state $s \in \mathcal{S}$ if and only if $\text{Pre}(a) \subseteq s$.
    \item The deterministic state transition function $\delta^*: \mathcal{S} \times \mathcal{A} \to \mathcal{S} \cup \{\bot\}$ is defined as:
    \begin{equation}
        \delta^*(s, a) = \begin{cases}
            (s \setminus \text{Del}^*(a)) \cup \text{Add}^*(a) & \text{if } \text{Pre}^*(a) \subseteq s \\
            \bot & \text{otherwise}
        \end{cases}
    \end{equation}
    where $\bot$ indicates execution failure (inapplicability).
\end{itemize}

A learned or candidate world model is denoted by $\widehat{M} = \langle \mathcal{F}, \mathcal{A}, \widehat{\delta}, s_0, S_g \rangle$, with candidate action schemas $\widehat{a} = \langle \widehat{\text{Pre}}(a), \widehat{\text{Add}}(a), \widehat{\text{Del}}(a) \rangle$.

\begin{definition}[Sparse Rule Intervention]
The syntactic distance between ground-truth model $M^*$ and candidate model $\widehat{M}$ is defined by the Action-Language AST edit distance:
\begin{equation}
    \Delta DSL(M^*, \widehat{M}) = \sum_{a \in \mathcal{A}} \left( |\text{Pre}^*(a) \triangle \widehat{\text{Pre}}(a)| + |\text{Add}^*(a) \triangle \widehat{\text{Add}}(a)| + |\text{Del}^*(a) \triangle \widehat{\text{Del}}(a)| \right) = k
\end{equation}
When $k \ll |\mathcal{A}| \cdot |\mathcal{F}|$, the intervention is termed \emph{sparse}. A \emph{precondition omission} occurs when $\widehat{\text{Pre}}(a) \subset \text{Pre}^*(a)$ while effect sets are unmodified ($\widehat{\text{Add}}(a) = \text{Add}^*(a)$, $\widehat{\text{Del}}(a) = \text{Del}^*(a)$).
\end{definition}

% ===========================================================================
\section{Formal Metric Definitions}
% ===========================================================================

\subsection{State-Indexed Search Trees}

\begin{definition}[Search Tree]
For search horizon depth $D \in \mathbb{N}$, let $G_T = (V_T, E_T)$ denote the state-indexed search tree generated by exhaustive forward search from initial state $s_0$ under ground-truth dynamics $M^*$.
Each node $u \in V_T$ corresponds to a reachable state-trajectory prefix $s(u) \in \mathcal{S}$. Each directed labeled edge $e = (u, a, v) \in E_T$ denotes an executable transition satisfying $\delta^*(s(u), a) = s(v) \ne \bot$.
Similarly, let $\widehat{G}_T = (\widehat{V}_T, \widehat{E}_T)$ denote the search tree generated by candidate model $\widehat{M}$ from $s_0$ up to depth $D$.
\end{definition}

\subsection{Local Transition Prediction Accuracy (\texorpdfstring{$A_{\text{pred}}$}{A_pred})}

Local evaluation assesses transition predictions over a finite dataset of transition queries $\mathcal{D}$.

\begin{definition}[Local Accuracy]
Given a test dataset $\mathcal{D} = \{ (s_i, a_i, s_i^*) \}_{i=1}^N$ where $s_i^* = \delta^*(s_i, a_i)$, the local prediction accuracy is:
\begin{equation}
    A_{\text{pred}}(\widehat{M}; \mathcal{D}) = \frac{1}{|\mathcal{D}|} \sum_{(s, a, s^*) \in \mathcal{D}} \mathbb{I}\left( \widehat{\delta}(s, a) = s^* \right)
\end{equation}
We formalize two canonical test distributions:
\begin{enumerate}
    \item \textbf{Reachable Valid Traces ($\mathcal{D}_{\text{valid}}$)}: Evaluated strictly along true environment transitions:
    \begin{equation}
        \mathcal{D}_{\text{valid}} = \{ (s(u), a, s(v)) \mid (u, a, v) \in E_T(M^*) \}
    \end{equation}
    \item \textbf{Search Fringe ($\mathcal{D}_{\text{fringe}}$)}: Evaluated across the union of physical transitions and candidate model expansions:
    \begin{equation}
        \mathcal{D}_{\text{fringe}} = E_T(M^*) \cup (\widehat{E}_T \setminus E_T(M^*))
    \end{equation}
\end{enumerate}
\end{definition}

\subsection{Global Execution Correctness: \texorpdfstring{$\text{PESR}$}{PESR} and \texorpdfstring{$R_{\text{play}}$}{R_play}}

Let $\operatorname{Planner}(\widehat{M}, s_0, g)$ denote an optimal cost-minimizing forward search planner (e.g., $A^*$ with an admissible heuristic $h^{\text{LM-cut}}$) operating on candidate model $\widehat{M}$. The planner synthesizes a candidate plan $\pi^*_{\widehat{M}} = \langle a_1, a_2, \dots, a_m \rangle$.

\begin{definition}[Plan Execution Success Rate ($\text{PESR}$)]
$\text{PESR}$ measures whether the plan synthesized on model $\widehat{M}$ executes successfully in the physical environment $M^*$ to reach goal $g \subseteq \mathcal{F}$:
\begin{equation}
    \text{PESR}(\widehat{M}) = \begin{cases}
        1 & \text{if } \delta^*(s_0, \pi^*_{\widehat{M}}) \models g \\
        0 & \text{otherwise}
    \end{cases}
\end{equation}
\end{definition}

\begin{definition}[Play Regret ($R_{\text{play}}$)]
Let $C_{M^*}(\pi)$ denote the cumulative execution cost of plan $\pi$ in the ground-truth environment $M^*$, defined as:
\begin{equation}
    C_{M^*}(\pi) = \begin{cases}
        \sum_{t=1}^{|\pi|} \operatorname{cost}(a_t) & \text{if } \forall t,\, \delta^*(s_{t-1}, a_t) \ne \bot \text{ and } \delta^*(s_0, \pi) \models g \\
        \infty & \text{if } \exists t,\, \delta^*(s_{t-1}, a_t) = \bot \text{ (execution failure)}
    \end{cases}
\end{equation}
The Play Regret of model $\widehat{M}$ relative to the ground-truth optimal plan $\pi^*_{M^*}$ is:
\begin{equation}
    R_{\text{play}}(\widehat{M}) = C_{M^*}(\pi^*_{\widehat{M}}) - C_{M^*}(\pi^*_{M^*})
\end{equation}
Whenever an inexecutable plan is synthesized, $R_{\text{play}}(\widehat{M}) = \infty$.
\end{definition}

\subsection{Topological Search-Tree Gap Metric (\texorpdfstring{$d_\triangle$}{d_triangle})}

\begin{definition}[Topological Edit Distance]
The structural gap between the physical search tree $G_T$ and the model search tree $\widehat{G}_T$ is defined as the symmetric difference of their edge sets:
\begin{equation}
    d_\triangle(G_T, \widehat{G}_T) = |E_T \triangle \widehat{E}_T| = |E_T \setminus \widehat{E}_T| + |\widehat{E}_T \setminus E_T|
\end{equation}
Edges in $E_{\text{phantom}} = \widehat{E}_T \setminus E_T$ are termed \emph{phantom edges}.
\end{definition}

% ===========================================================================
\section{Formal Proofs}
% ===========================================================================

\subsection{Lemma 1: Monotone Edge Addition under Precondition Omission}

\begin{lemma}[No-Edges-Lost Lemma]\label{lem:no_edges_lost}
Let $M^* = \langle \mathcal{F}, \mathcal{A}, \delta^*, s_0, S_g \rangle$ be the ground-truth transition system. Suppose $\widehat{M}$ is constructed such that for all actions $a \in \mathcal{A}$:
\begin{equation}
    \widehat{\text{Pre}}(a) \subseteq \text{Pre}^*(a), \quad \widehat{\text{Add}}(a) = \text{Add}^*(a), \quad \widehat{\text{Del}}(a) = \text{Del}^*(a)
\end{equation}
Then every edge in the true search tree $G_T$ is preserved in candidate search tree $\widehat{G}_T$:
\begin{equation}
    E_T(M^*) \subseteq \widehat{E}_T \quad \text{and} \quad |E_T \setminus \widehat{E}_T| = 0
\end{equation}
Consequently, the topological gap simplifies strictly to the number of phantom edges:
\begin{equation}
    d_\triangle(G_T, \widehat{G}_T) = |\widehat{E}_T \setminus E_T| \ge 0
\end{equation}
\end{lemma}

\begin{proof}
Let $e = (u, a, v) \in E_T(M^*)$ be an arbitrary edge in the ground-truth search tree $G_T$.
By definition of $G_T$, this edge corresponds to a valid physical transition:
\begin{equation}
    \delta^*(s(u), a) = s(v) \ne \bot
\end{equation}
This implies that action $a$ is applicable at state $s(u)$ in $M^*$:
\begin{equation}
    \text{Pre}^*(a) \subseteq s(u)
\end{equation}
By hypothesis, the candidate model satisfies $\widehat{\text{Pre}}(a) \subseteq \text{Pre}^*(a)$. By transitivity of subset inclusion:
\begin{equation}
    \widehat{\text{Pre}}(a) \subseteq s(u)
\end{equation}
Hence, action $a$ is applicable at state $s(u)$ in $\widehat{M}$.
Next, compute the successor state $\widehat{\delta}(s(u), a)$:
\begin{equation}
    \widehat{\delta}(s(u), a) = (s(u) \setminus \widehat{\text{Del}}(a)) \cup \widehat{\text{Add}}(a)
\end{equation}
Because effect sets are unmodified ($\widehat{\text{Del}}(a) = \text{Del}^*(a)$ and $\widehat{\text{Add}}(a) = \text{Add}^*(a)$):
\begin{equation}
    \widehat{\delta}(s(u), a) = (s(u) \setminus \text{Del}^*(a)) \cup \text{Add}^*(a) = \delta^*(s(u), a) = s(v)
\end{equation}
We now proceed by induction on search tree depth $t \in [0, D]$.
\begin{itemize}
    \item \textbf{Base Case ($t=0$)}: At the root node $u_0$, $s(u_0) = s_0 \in V_T$ and $s_0 \in \widehat{V}_T$.
    \item \textbf{Inductive Step}: Assume that for all nodes $u$ at depth $t < D$, $u \in V_T \implies u \in \widehat{V}_T$. For any edge $(u, a, v) \in E_T$, action $a$ is applicable in $\widehat{M}$ and produces identical successor $s(v)$. Thus node $v \in \widehat{V}_T$ and edge $(u, a, v) \in \widehat{E}_T$.
\end{itemize}
Therefore, $E_T \subseteq \widehat{E}_T$, which implies $E_T \setminus \widehat{E}_T = \emptyset$, so $|E_T \setminus \widehat{E}_T| = 0$.
Substituting into the definition of symmetric difference:
\begin{equation}
    d_\triangle(G_T, \widehat{G}_T) = |E_T \setminus \widehat{E}_T| + |\widehat{E}_T \setminus E_T| = 0 + |\widehat{E}_T \setminus E_T| = |\widehat{E}_T \setminus E_T|
\end{equation}
This completes the proof.
\end{proof}

% ---------------------------------------------------------------------------
\subsection{Proposition 1: Search-Tree Phantom Path Divergence}

\begin{proposition}[Topological Divergence Bound]\label{prop:phantom_divergence}
Let $G_T$ be a regular $b$-ary search tree of uniform branching factor $b \ge 2$ and depth horizon $D \in \mathbb{N}$. Suppose candidate model $\widehat{M}$ is constructed by omitting a single bottleneck precondition $p^* \in \text{Pre}^*(a^*)$ for an action $a^*$ applicable at depth $d < D$.
Under the Fresh-Subtree Assumption (the unblocked state expands a regular tree of depth $D - (d + 1)$ with branching factor $b$), the topological gap and local prediction accuracy satisfy:
\begin{equation}
    d_\triangle(G_T, \widehat{G}_T) = \Omega\left(b^{D-d}\right)
\end{equation}
and evaluated over the candidate search distribution $\mathcal{D}_{\text{fringe}} = E_T \cup (\widehat{E}_T \setminus E_T)$:
\begin{equation}
    A_{\text{pred}}(\widehat{M}; \mathcal{D}_{\text{fringe}}) \ge 1 - b^{-d}
\end{equation}
Furthermore, for any precision threshold $\epsilon > 0$, whenever the bottleneck depth satisfies $d \ge \lceil \log_b(1/\epsilon) \rceil$:
\begin{equation}
    A_{\text{pred}}(\widehat{M}) \ge 1 - \epsilon \quad \text{while} \quad d_\triangle(G_T, \widehat{G}_T) \xrightarrow{D \to \infty} \infty
\end{equation}
\end{proposition}

\begin{proof}
We count the cardinality of the respective edge sets.
In a regular $b$-ary search tree of uniform depth $D$, the total number of physical edges in $G_T$ is:
\begin{equation}
    |E_T| = \sum_{i=0}^{D-1} b^{i+1} = b \cdot \frac{b^D - 1}{b - 1} = \Theta(b^D)
\end{equation}
At depth $d \in [0, D-1]$, omitting precondition $p^*$ unblocks action $a^*$ at a state $s_d$, creating an initial phantom edge $(s_d, a^*, s_{\text{phantom}})$.
By the Fresh-Subtree Assumption, the node $s_{\text{phantom}}$ is expanded recursively to horizon $D$, forming a complete subtree of depth $D - (d + 1)$. The number of phantom edges generated in this unblocked subtree is:
\begin{equation}
    |E_{\text{phantom}}| = 1 + \sum_{j=1}^{D-d-1} b^j = 1 + b \cdot \frac{b^{D-d-1} - 1}{b - 1} = \frac{b^{D-d} - 1}{b - 1} = \Theta\left(b^{D-d}\right)
\end{equation}
By \Cref{lem:no_edges_lost}, $E_T \subseteq \widehat{E}_T$, so the topological gap is:
\begin{equation}
    d_\triangle(G_T, \widehat{G}_T) = |E_{\text{phantom}}| = \frac{b^{D-d} - 1}{b - 1} = \Omega\left(b^{D-d}\right)
\end{equation}
Now consider local prediction accuracy over the candidate search distribution:
\begin{equation}
    \mathcal{D}_{\text{fringe}} = E_T \cup E_{\text{phantom}}, \quad |\mathcal{D}_{\text{fringe}}| = |E_T| + |E_{\text{phantom}}|
\end{equation}
By \Cref{lem:no_edges_lost}, for every edge in $E_T$, model $\widehat{M}$ predicts the exact ground-truth transition. The model errs exclusively on the phantom transitions in $E_{\text{phantom}}$, where $\widehat{M}$ predicts a valid state change while ground truth $M^*$ returns $\bot$.
Thus:
\begin{equation}
    A_{\text{pred}}(\widehat{M}; \mathcal{D}_{\text{fringe}}) = \frac{|E_T|}{|E_T| + |E_{\text{phantom}}|}
\end{equation}
Substituting the exact sums:
\begin{align}
    A_{\text{pred}}(\widehat{M}) &= \frac{b \frac{b^D - 1}{b - 1}}{b \frac{b^D - 1}{b - 1} + \frac{b^{D-d} - 1}{b - 1}} = \frac{b(b^D - 1)}{b(b^D - 1) + (b^{D-d} - 1)} \\
    &= \frac{1}{1 + \frac{b^{D-d} - 1}{b(b^D - 1)}} = \frac{1}{1 + b^{-d} \left( \frac{1 - b^{-(D-d)}}{b(1 - b^{-D})} \right)}
\end{align}
Since $b \ge 2$ and $D > d \ge 1$, we have:
\begin{equation}
    \frac{1 - b^{-(D-d)}}{b(1 - b^{-D})} \le \frac{1}{b} < 1
\end{equation}
Using the inequality $\frac{1}{1 + x} \ge 1 - x$ for $x \ge 0$:
\begin{equation}
    A_{\text{pred}}(\widehat{M}; \mathcal{D}_{\text{fringe}}) \ge 1 - b^{-d}
\end{equation}
Finally, choose any $\epsilon > 0$. Setting $d \ge \lceil \log_b(1/\epsilon) \rceil$ implies $b^{-d} \le \epsilon$, guaranteeing:
\begin{equation}
    A_{\text{pred}}(\widehat{M}; \mathcal{D}_{\text{fringe}}) \ge 1 - \epsilon
\end{equation}
Simultaneously, holding $d$ fixed and taking the horizon limit $D \to \infty$:
\begin{equation}
    \lim_{D \to \infty} d_\triangle(G_T, \widehat{G}_T) = \lim_{D \to \infty} \frac{b^{D-d} - 1}{b - 1} = \infty
\end{equation}
This concludes the proof.
\end{proof}

% ---------------------------------------------------------------------------
\subsection{Proposition 2: Causal Decoupling and Statistical Independence of \texorpdfstring{$A_{\text{pred}}$}{A_pred} and \texorpdfstring{$\text{PESR}$}{PESR}}

\begin{proposition}[Causal Decoupling Theorem]\label{prop:independence}
Local prediction accuracy $A_{\text{pred}}$ and global plan execution success rate $\text{PESR}$ are functionally and statistically decoupled. Specifically:
\begin{enumerate}
    \item \textbf{High Accuracy with Total Execution Collapse}: For any arbitrary tolerance $\epsilon > 0$, there exists an infinite family of planning tasks $\Pi$ and candidate models $\widehat{M}$ such that:
    \begin{equation}
        A_{\text{pred}}(\widehat{M}) \ge 1 - \epsilon \quad \text{and} \quad \text{PESR}(\widehat{M}) = 0 \quad (R_{\text{play}}(\widehat{M}) = \infty)
    \end{equation}
    \item \textbf{Low Accuracy with Perfect Optimal Execution}: Conversely, for any $\epsilon > 0$, there exist models $\widehat{M}'$ exhibiting arbitrarily low local accuracy:
    \begin{equation}
        A_{\text{pred}}(\widehat{M}') \le \epsilon \quad \text{and} \quad \text{PESR}(\widehat{M}') = 1 \quad (R_{\text{play}}(\widehat{M}') = 0)
    \end{equation}
\end{enumerate}
Consequently, there exists no non-trivial monotonic lower-bounding function $f: [0, 1] \to [0, 1]$ satisfying $\text{PESR}(\widehat{M}) \ge f(A_{\text{pred}}(\widehat{M}))$ with $f(1-\epsilon) > 0$.
\end{proposition}

\begin{proof}
We prove both parts constructively.

\paragraph{Part 1: Constructing $A_{\text{pred}} \ge 1 - \epsilon$ with $\text{PESR} = 0$.}
Construct a parameterized planning task $\Pi_d$ with horizon $D \gg d$.
Let the physical environment $M^*$ possess an optimal execution path $\pi^* = \langle a_1^*, a_2^*, \dots, a_L^* \rangle$ from $s_0$ to goal $g$ with path length $L = D$ and uniform cost $C(\pi^*) = L$.
At search depth $d = \lceil \log_b(1/\epsilon) \rceil$, introduce an action $a_{\text{shortcut}}$ that is physically blocked by bottleneck precondition $p^* \in \text{Pre}^*(a_{\text{shortcut}})$ (i.e., $s_d \not\models p^*$).
If $a_{\text{shortcut}}$ were executable, its transition reaches goal state $g$ directly with action cost $1$.
Thus, any path utilizing $a_{\text{shortcut}}$ has total planned cost:
\begin{equation}
    C(\pi_{\text{phantom}}) = d + 1 < L
\end{equation}
Now, let candidate model $\widehat{M}$ be constructed by omitting $p^*$ from $a_{\text{shortcut}}$ ($\widehat{\text{Pre}}(a_{\text{shortcut}}) = \text{Pre}^*(a_{\text{shortcut}}) \setminus \{p^*\}$), while all other actions and effects are perfectly specified.
\begin{itemize}
    \item By \Cref{prop:phantom_divergence}, the local accuracy of $\widehat{M}$ satisfies:
    \begin{equation}
        A_{\text{pred}}(\widehat{M}; \mathcal{D}_{\text{fringe}}) \ge 1 - b^{-d} \ge 1 - \epsilon
    \end{equation}
    \item We now evaluate the cost-minimizing planner $\operatorname{Planner}(\widehat{M}, s_0, g)$.
    Because the planner searches for minimal-cost paths, it evaluates:
    \begin{equation}
        \operatorname{cost}(\pi_{\text{phantom}}) = d + 1 < L = \operatorname{cost}(\pi^*)
    \end{equation}
    Therefore, the planner deterministically selects plan $\pi^*_{\widehat{M}} = \pi_{\text{phantom}} = \langle a_1^*, \dots, a_d^*, a_{\text{shortcut}} \rangle$.
    \item When $\pi^*_{\widehat{M}}$ is executed in the physical environment $M^*$, prefix $\langle a_1^*, \dots, a_d^* \rangle$ executes successfully to reach state $s_d$.
    At step $d+1$, action $a_{\text{shortcut}}$ is executed. Since $s_d \not\models p^*$, we have:
    \begin{equation}
        \delta^*(s_d, a_{\text{shortcut}}) = \bot
    \end{equation}
    The plan halts catastrophically before reaching goal $g$.
    Consequently, $\text{PESR}(\widehat{M}) = 0$ and $R_{\text{play}}(\widehat{M}) = \infty$.
\end{itemize}

\paragraph{Part 2: Constructing $A_{\text{pred}} \le \epsilon$ with $\text{PESR} = 1$.}
Let $\pi^* = \langle a_1^*, a_2^*, \dots, a_m^* \rangle$ be the ground-truth optimal plan for task $\Pi$, traversing states $\langle s_0, s_1^*, \dots, s_m^* \rangle$ with $s_m^* \models g$.
Augment the environment with a set of $K$ distractor actions $\mathcal{A}_{\text{distractor}}$ that govern task-irrelevant fluents $\mathcal{F}_{\text{distractor}}$, where $|\mathcal{A}_{\text{distractor}}| \ge \frac{m}{\epsilon}$.
Construct model $\widehat{M}'$ such that:
\begin{enumerate}
    \item $\widehat{M}'$ specifies the actions along the optimal path $\{a_1^*, \dots, a_m^*\}$ with $100\%$ fidelity.
    \item For all distractor actions $a \in \mathcal{A}_{\text{distractor}}$, $\widehat{M}'$ predicts completely incorrect transitions (e.g., predicting $\bot$ for valid transitions or invalid state updates).
\end{enumerate}
Evaluating local accuracy over the full transition distribution yields:
\begin{equation}
    A_{\text{pred}}(\widehat{M}') \le \frac{m}{m + |\mathcal{A}_{\text{distractor}}|} \le \frac{m}{m + m/\epsilon} \le \epsilon
\end{equation}
However, when the forward planner operates on $\widehat{M}'$, all distractor actions are either perceived as inexecutable or leading to unpromising states. The planner deterministically outputs the unique valid path $\pi^*$.
Executing $\pi^*$ in $M^*$ succeeds with zero regret:
\begin{equation}
    \text{PESR}(\widehat{M}') = 1 \quad \text{and} \quad R_{\text{play}}(\widehat{M}') = C_{M^*}(\pi^*) - C_{M^*}(\pi^*) = 0
\end{equation}

\paragraph{Non-Existence of Monotone Bound.}
Suppose for contradiction there exists a function $f: [0, 1] \to [0, 1]$ such that $\text{PESR}(\widehat{M}) \ge f(A_{\text{pred}}(\widehat{M}))$ for all $\widehat{M}$, with $f(1-\epsilon) > 0$.
Applying this bound to Part 1 implies:
\begin{equation}
    0 = \text{PESR}(\widehat{M}) \ge f(A_{\text{pred}}(\widehat{M})) \ge f(1 - \epsilon) > 0
\end{equation}
which is an immediate contradiction. Hence, no such lower bound exists, establishing that local prediction accuracy and global plan execution are functionally decoupled.
\end{proof}

% ===========================================================================
\section{Conclusion}
% ===========================================================================

The formal proofs established in \Cref{lem:no_edges_lost}, \Cref{prop:phantom_divergence}, and \Cref{prop:independence} provide the theoretical foundation for this thesis. They prove that the verified-vs-correct gap in symbolic world models is not an empirical artifact or a quirk of specific training runs, but a fundamental property of combinatorial search over state-space graphs. Evaluating world models via passive transition accuracy is mathematically incapable of certifying autonomous agents for goal-directed play. Resolving this challenge necessitates active, counterexample-guided repair mechanisms such as CEG-OMR.

\end{document}
